\chapter{Cortisol}

\section*{Definition and Overview}
Cortisol is a steroid hormone produced by the adrenal glands, two small glands sitting on top of the kidneys. It is often called the stress hormone because it is released in response to physical and emotional stress, but its role is far broader: cortisol helps regulate blood sugar, blood pressure, the immune system, and the sleep--wake cycle. Cortisol levels follow a daily rhythm, peaking in the early morning and falling overnight. Understanding cortisol matters because both too little and too much of it cause illness -- Addison's disease when levels are low, and Cushing's syndrome when they are high.

\section*{Epidemiology}
Disorders of cortisol are uncommon. Cushing's syndrome, caused by excess cortisol, affects roughly 1--2 people per 100,000 per year, most often women aged 20--50, and is frequently due to prolonged use of steroid medicines. Addison's disease, or primary adrenal insufficiency, is rare, affecting about 1 in 10,000 people. Mild disturbances of the stress response are widely discussed in popular health culture, but clinically significant cortisol abnormalities are unusual. When they occur, they are serious and require specialist investigation and management.

\section*{Aetiology and Pathophysiology}
Cortisol production is controlled by a feedback loop involving the brain and the adrenal glands.

\begin{itemize}
    \item \textbf{Production:} the hypothalamus releases CRH, prompting the pituitary to release ACTH, which stimulates the adrenal cortex to make cortisol
    \item \textbf{Stress response:} physical illness, injury, surgery, and psychological stress increase cortisol release
    \item \textbf{Metabolic effects:} cortisol raises blood sugar, mobilises fat and protein, and supports energy in stress
    \item \textbf{Immune effects:} cortisol dampens inflammation, which is why steroids treat allergic and inflammatory disease
    \item \textbf{Circulating rhythm:} levels are highest around waking and lowest in the early night
    \item \textbf{Excess cortisol:} from steroid medicines, adrenal tumours, or pituitary tumours (Cushing's syndrome)
    \item \textbf{Cortisol deficiency:} from adrenal gland destruction (Addison's disease) or pituitary failure
\end{itemize}

\section*{Clinical Features}
Symptoms depend on whether cortisol is too high or too low.

\begin{itemize}
    \item \textbf{High cortisol (Cushing's syndrome):} weight gain, especially around the face, abdomen, and neck; thinning skin, easy bruising, and purple stretch marks; high blood pressure and blood sugar; muscle weakness; mood changes; irregular periods
    \item \textbf{Low cortisol (Addison's disease):} fatigue, weakness, weight loss, low blood pressure and dizziness on standing, salt craving, nausea, and darkening of the skin
    \item \textbf{Normal variations:} stress-related tiredness and sleep disturbance without a true cortisol disorder
\end{itemize}

\section*{Differential Diagnosis}
Symptoms of cortisol disorders overlap with many common conditions.

\begin{itemize}
    \item Chronic fatigue syndrome, depression, and anxiety
    \item Obesity, metabolic syndrome, and type 2 diabetes (for Cushing's symptoms)
    \item Hypothyroidism, with fatigue and weight change
    \item Anaemia and other causes of fatigue and dizziness
    \item Polycystic ovary syndrome, for menstrual irregularity and weight gain
    \item Malignancy or chronic infection causing weight loss
    \item Conditions causing low blood pressure, such as dehydration or heart disease
\end{itemize}

\section*{When to Seek Medical Care}
See a doctor if you have persistent, unexplained fatigue, weight change, high blood pressure, easy bruising, or other symptoms that might suggest a cortisol disorder. Tests are simple and can clarify whether cortisol levels are normal.

\textbf{Call 000 (or local emergency services) for:}
\begin{itemize}
    \item An adrenal (addisonian) crisis: sudden severe weakness, vomiting, confusion, collapse, or very low blood pressure in someone known to have adrenal insufficiency
    \item Severe abdominal or back pain with collapse
    \item A severe infection or illness in a person taking long-term steroids, who may need extra steroid cover
    \item Sudden neurological symptoms or chest pain
\end{itemize}

\section*{Diagnosis and Investigations}
Diagnosis requires specific hormonal testing, interpreted by an endocrinologist.

\begin{itemize}
    \item \textbf{Morning cortisol blood test:} a low early-morning cortisol suggests adrenal insufficiency; a high value can suggest excess
    \item \textbf{ACTH measurement:} helps distinguish adrenal from pituitary causes
    \item \textbf{Salivary or urinary cortisol:} measured over 24 hours or at night to assess total cortisol production
    \item \textbf{Dexamethasone suppression test:} a low-dose steroid is given and cortisol measured, used to confirm Cushing's syndrome
    \item \textbf{ACTH stimulation test:} synthetic ACTH is given to test the adrenal response
    \item \textbf{Imaging:} CT of the adrenal glands or MRI of the pituitary to find tumours
    \item \textbf{Other blood tests:} electrolytes, blood sugar, and thyroid function
\end{itemize}

\section*{Management}
Treatment depends on whether cortisol is excessive or deficient.

\begin{itemize}
    \item \textbf{Adrenal insufficiency:} lifelong replacement with hydrocortisone (or prednisolone) and fludrocortisone, with dose increases during illness ("sick day rules")
    \item \textbf{Adrenal crisis:} emergency treatment with intravenous hydrocortisone, fluids, and glucose
    \item \textbf{Cushing's syndrome from steroids:} gradual reduction of steroid medicines under medical supervision, never stopping abruptly
    \item \textbf{Cushing's from a tumour:} surgical removal of the adrenal or pituitary tumour, with radiotherapy or medicines as alternatives
    \item \textbf{Medicines:} drugs that block cortisol production for inoperable tumours
    \item \textbf{Monitoring:} regular endocrinology review, blood pressure control, and bone health
    \item \textbf{Education:} patients with adrenal insufficiency carry a warning card or bracelet and learn sick-day rules
\end{itemize}

\section*{Complications}
\begin{itemize}
    \item Adrenal crisis, a life-threatening emergency in untreated or poorly managed adrenal insufficiency
    \item Complications of Cushing's syndrome: diabetes, hypertension, osteoporosis, infections, and blood clots
    \item Poor wound healing and skin fragility from high cortisol
    \item Mood disturbance and depression in both excess and deficiency
    \item Complications of long-term steroid treatment, including weight gain and bone loss
    \item Complications of pituitary or adrenal surgery
    \item Growth failure in children with excess cortisol
\end{itemize}

\section*{Prognosis and Outcomes}
With correct diagnosis and treatment, people with cortisol disorders generally do very well. Adrenal insufficiency is managed safely with replacement hormones, and most people live full, active lives, provided they follow sick-day rules during illness. Cushing's syndrome usually improves markedly after removal of the cause, though blood pressure, diabetes, and bone density may take time to recover and can need ongoing treatment. Untreated, both excess and deficiency of cortisol are serious and can be fatal, which is why these conditions should never be managed without specialist care.

\section*{Prevention and Self-Care}
\begin{itemize}
    \item Do not stop long-term steroid medicines abruptly; always taper under medical supervision
    \item If you have adrenal insufficiency, carry identification and know your sick-day rules
    \item Double your steroid dose during fever, infection, surgery, or significant illness, as instructed by your doctor
    \item Keep a supply of injectable hydrocortisone if prescribed
    \item Attend endocrinology follow-up and have regular blood tests
    \item Manage stress, sleep, and lifestyle to support general wellbeing
    \item Seek medical review for persistent unexplained symptoms rather than self-diagnosing
\end{itemize}

\section*{Sources and Further Reading}
The following reputable sources provide further information for healthcare students and patients seeking reliable, current advice about cortisol and adrenal disorders.

\begin{itemize}
    \item NHS. \textit{Addison's disease and adrenal insufficiency.} https://www.nhs.uk/conditions/addisons-disease/
    \item Mayo Clinic. \textit{Cushing syndrome.} https://www.mayoclinic.org/diseases-conditions/cushing-syndrome/
    \item MSD Manual. \textit{Adrenal gland disorders: professional reference.} https://www.msdmanuals.com/
    \item MedlinePlus. \textit{Cortisol test and adrenal disorders.} https://medlineplus.gov/
    \item Endocrine Society. \textit{Patient resources on adrenal disorders.} https://www.endocrine.org/
    \item Addison's Disease Self-Help Group Australia. \textit{Patient support and education.} https://www.addisons.org.au/
\end{itemize}
